\chapter{数列（二）}
\section{数列的通项公式}
\subsection{累加与累乘}
\clearpage
\subsection{待定系数法的辅助数列}
\clearpage
\subsection{不动点与特征根}
\clearpage
\subsection{分奇偶讨论}
\begin{example}
    数列\(\{a_n\}\) 满足 \(a_1 = 2\)，\(a_{n+1} + a_n = 4n - 3\)，求 \(\{a_n\}\) 的通项公式。
\end{example}
\begin{exercise}
    数列\(\{a_n\}\) 满足 \(a_1 = 1\)，\(a_{n+1} + a_n = 2^{n-1}\)，求 \(\{a_n\}\) 的通项公式。
\end{exercise}

奇偶项数列能求通项的条件是两侧下标奇偶性相同。

\begin{example}
数列 \(\{a_n\}\) 满足 \(a_1 = 3\)，\(a_n a_{n+1} = 9 \times 2^{2n-1}\)，求数列 \(\{a_n\}\) 的通项公式。
\end{example}

\begin{example}
    数列 \(\{a_n\}\) 满足 \(a_1=1, a_2 = 2\)，\(a_{n+2} + (-1)^{n+1} a_n = 1 + (-1)^n\)，求数列 \(\{a_n\}\) 的通项公式。
\end{example}

\begin{example}
    \(\{a_n\}\) 满足 \(a_1 = 1\)，\(a_{n+1} = \begin{cases} 
        a_n + 1 & n \text{为奇} \\
        2a_n & n \text{为偶}
        \end{cases}\)，求 \(\{a_n\}\) 的通项公式。
\end{example}
\begin{exercise}
    已知数列 \(\{a_n\}\) 满足 \(a_1 = 1\)，\(a_{n+1} = 
\begin{cases} 
a_n + 1, & n \text{为奇数}, \\
a_n + 2, & n \text{为偶数}.
\end{cases}\)，求数列 \(\{a_n\}\) 的通项公式．

\end{exercise}

\begin{example}
    在数列 \(\{a_n\}\) 中, \(a_1 = 0\)，且对任意 \(k \in \mathbb{N}^*\)，\(a_{2k-1}\)，\(a_{2k}\)，\(a_{2k+1}\) 成等差数列，其公差为 \(d_k\)。

\begin{enumerate}
    \item 若 \(d_k = 2k\)，求 \(\{a_n\}\) 的通项公式并证明 \(a_{2k}\)，\(a_{2k+1}\)，\(a_{2k+2}\) 成等比数列，\(k \in \mathbb{N}^*\)；
    \item 若对任意 \(k \in \mathbb{N}^*\)，\(a_{2k}\)，\(a_{2k+1}\)，\(a_{2k+2}\) 成等比数列，其公比为 \(q_k\)，设 \(q_1 \neq 1\)，证明 \(\left\{\dfrac{1}{q_k - 1}\right\}\) 成等差数列。
\end{enumerate}
\end{example}

\clearpage
